\cotefile{58}
\scanpage{20}
\mspage{- 2 -}
\noindent\hfill n\textsuperscript{o}358

$V'_{(K)}$ dans $W$ et une seule telle que $\bar f'(v'\handwritten{\otimes} 1)=f'(v')$ pour tout $v'$ dans $V'$. Si $\handwritten{\lambda}$ est un isomorphisme, l'application $K$-linéaire $\bar f':\handwritten{\lambda}^{-1}$ est alors l'unique application $K$-linéaire de $V$ dans $W$ prolongeant $f'$ ; autrement dit, a) implique b). D'autre part, si b) est satisfaite, il existe une application $K$-linéaire $\mu$ de $V$ dans $V'_{(K)}$ telle que $\mu(v')=v'\handwritten{\otimes} 1$ pour tout $v'$ dans $V'$ ; il est immédiat que $\handwritten{\lambda}$ et $\mu$ sont des applications réciproques, et par suite $\handwritten{\lambda}$ est une bijection. On a donc prouvé que a) équivaut à b).

Par ailleurs, pour toute base $B'$ de $V'$ sur $K'$, l'ensemble $B$ des éléments $v'\handwritten{\otimes} 1$ de $V'_{(K)}$ (pour $v'$ parcourant $B'$) est une base de $V'_{\handwritten{(K)}}$ sur $K$, et $\handwritten{\lambda}$ transforme $B$ en $B'$. Pour que l'application $K$-linéaire $\handwritten{\lambda}$ soit bijective, il faut que l'image par $\handwritten{\lambda}$ de toute base de $V'_{(K)}$ soit une base de $V$ sur $K$, et il suffit qu'il en soit ainsi pour une seule base. Ceci prouve l'équivalence de a), c) et d).

Comme toute partie de $V'$ libre sur $K'$ est contenue dans une base de $V'$ sur $K'$, c) implique e). Enfin, supposons e) satisfaite ; si $B'$ est une base de $V'$ sur $K'$, c'est une partie de $V$ libre sur $K$, et comme $B'$ engendre $V'$ sur $K'$, donc $V$ sur $K$ d'après e), on voit que $B'$ est une base de $V$ sur $K$ ; on a donc prouvé que e) implique c). C.Q.F.D.

DEFINITION 1.- \uline{Soient} $V$ \uline{un espace vectoriel à droite sur un corps} $K$ \uline{et} $K'$ \uline{un sous-corps de} $K$. \uline{On appelle} $K'$-\uline{structure sur} $V$ \uline{tout sous-espace} $K'$-\uline{vectoriel} $V'$ \uline{de} $V$ \uline{qui vérifie les conditions équivalentes de la prop}.1. \uline{On dit que les éléments de} $V'$ \uline{sont les vecteurs rationnels sur} $K'$.

2. \uline{Rationalité pour un sous-espace}.

DEFINITION 2.- \uline{Soient} $V$ \uline{un espace vectoriel à droite sur le corps} $K$ \uline{et} $V'$ \uline{une} $K'$-\uline{structure sur} $V$. \uline{On dit qu'un sous-espace} $K$-\uline{vectoriel} $W$ \uline{de} $V$ \uline{est rationnel sur} $K'$ \uline{s'il est engendré sur} $K$ \uline{par} $W\handwritten{\cap} V'$.

% archive sheet number: 19
\scanpage{21}

\tikz[remember picture,overlay]{\coordinate (lgma);}\uline{Lemme 4} Soit $X$ un préschéma noethérien connexe, \inserted{dont \unsure{les composantes irréductibles|tte composante irréductible} de dim $\geq k$, et} tel que pour tt pt \uline{fermé} $x\in X$, $(\operatorname{Spec}\underline{\mathcal O}_{X,x})'$ soit connexe en dim $\geq k$\deleted{\unsure{en|et|ens}}. Alors $X$ est connexe en dim $\geq k$.\tikz[remember picture,overlay]{\coordinate (lgmb); \draw ([xshift=-5pt,yshift=-4pt]lgma) -- ([xshift=-5pt]lgma |- lgmb);}
% slanted
\margin{\uline{peut se remplacer par} : pour tt $x$ tel que $\dim\bar x<k$, $(\operatorname{Spec}\underline{\mathcal O}_{X,x})'$ connexe. Alors il suffit de renvoyer au lemme plus général dans EGA $\overline{\underline{\mathrm{IV}}}$ \S{} du \unsure{crit. jacob.|crit. local}, dernier N\textsuperscript{o}.}

\deleted{Supposons} \inserted{Soit} en effet \deleted{\unsure{pour|que}} $Z$ une partie fermée de $X$ de dim $<k$, prouvons que $Z$ ne disconnecte pas $X$. \deleted{\unsure{Pour le voir|Par le raisont|Par la condition}} S'il \unsure{le|la} disconnectait, Alors on aurait $U=X-Z=\deleted{X}U'\cup\deleted{X}U''$, avec $U',U''$ fermés dans $U$, $U'\cap U''=\varnothing$, Soient $\deleted{U}X'=\overline{U'}$, $\deleted{U}X''=\overline{U''}$, donc comme $U$ dense
\deleted{$X=(\underbrace{Z\cup X'})\cup X''$, \deleted{donc} donc $(Z\cup X')\cap X''\ne\varnothing$}
dans $X$, $X^{\deleted{\unsure{z|2}}}=\deleted{U}X'\cup\deleted{U}X''$, donc $X'\cap X''\ne\varnothing$.
Soit $z$ un pt fermé de $X'\cap X''$. Alors \unsure{s'introduisant|s'induisant sur} \unsure{un} $\operatorname{Spec}\underline{\mathcal O}_{X,z}=X_z$, on voit que $Z_z$ disconnecte $X_z$, donc $Z'_z$ disconnecte $X'_z$, absurde \unsure{....|- - - -}

% archive sheet number: 20
\scanpage{22}
\mspage{- 3 -}
\noindent\hfill n\textsuperscript{o}355

Le $B$-module $E\handwritten{\otimes}_A F$ est réunion des sous-modules $E'\handwritten{\otimes}_A F$, $E'$ parcourant l'ensemble des sous-modules de type fini de $E$. Si donc $\underline P$ appartient à $\operatorname{Ass}_B(E\handwritten{\otimes}_A F)$, il existe un $E'$ de type fini tel que $\underline P\handwritten{\in}\operatorname{Ass}_B(E'\handwritten{\otimes}_A F)$. D'aprè\unsure{s} (ii), il existe $\underline p\handwritten{\in}\operatorname{Ass}_A(E')$ tel que $\underline P\handwritten{\in}\operatorname{Ass}_B(F\handwritten{/}\underline p F)$ ; comme $\operatorname{Ass}_A(E')\handwritten{\subset}\operatorname{Ass}_A(E)$, cela achève de démontrer la proposition.

COROLLAIRE.- \uline{Si} $A$ \uline{est noethérien}, \uline{et si} $\underline P\handwritten{\in}\operatorname{Ass}_B(E\handwritten{\otimes}_A F)$, \uline{on a} $\handwritten{\pi}^{-1}(\underline P)\handwritten{\in}\operatorname{Ass}_A(E)$.

Cela résulte des propositions 1 et 2.

\uline{Remarques}.-

1) Supposons $A$ \uline{et} $B$ \uline{noethériens}. Soit $\underline P$ un idéal premier de $B$, et soit $\underline p=\handwritten{\pi}^{-1}(\underline P)$. Soit $S=A-\underline p$, et soit $k(\underline p)=S^{-1}(A\handwritten{/}\underline p)$ le corps des fractions de $A\handwritten{/}\underline p$. Puisque $\underline P$ contient $\handwritten{\pi}(\underline p)$, il définit un idéal premier de $B\handwritten{/}\underline p B$, et puisqu'il ne rencontre pas $S$, il définit aussi un idéal premier, noté $\underline P'$, de $S^{-1}(B\handwritten{/}\underline p B)=k(\underline p)\handwritten{\otimes}_A B$. Alors, \uline{pour que} $\underline P\handwritten{\in}\operatorname{Ass}_B(E\handwritten{\otimes}_A F)$, \uline{il faut et il suffit que} $\underline p\handwritten{\in}\operatorname{Ass}_A(E)$ \uline{et} $\underline P'\handwritten{\in}\operatorname{Ass}_{k(\underline p)\handwritten{\otimes}_A B}(k(\underline p)\handwritten{\otimes}_A F)$.

Vu la prop.2, tout revient à voir que les conditions

\[
\text{`` }\underline P\handwritten{\in}\operatorname{Ass}_B(F\handwritten{/}\underline p F)\text{ '' \quad et \quad `` }\underline P'\handwritten{\in}\operatorname{Ass}_{k(\underline p)\handwritten{\otimes}_A B}(k(\underline p)\handwritten{\otimes}_A F)\text{ ''}
\]

sont équivalentes, ce qui résulte de la prop.5 du \handwritten{\S} 1 (c'est ici que l'on utilise l'hypothèse noethérienne sur $B$).

\handwritten{[}D'après Grothendieck, l'énoncé ci-dessus est celui que l'on utilise dans les schémas (le produit tensoriel par $k(\underline p)$ s'interprétant comme le passage à la ``fibre au-dessus de $\underline p$'').\handwritten{]}

2) La proposition 2 s'applique notamment lorsque $B$ \uline{est plat sur} $A$ et que l'on prend $F=B$. Le module $E\handwritten{\otimes}_A F$ est alors simplement le module $E_{(B)}$ \uline{déduit de} $E$ \uline{par extension des scalaires}.

On recopie ensuite les deux dernières pages de l'Etat 1.

\begin{center}-:-:-:-\end{center}

% archive sheet number: 21
\scanpage{23}

\tikz[remember picture,overlay]{\coordinate (lgta);}\uline{Théorème 2} Soit $f:X\longrightarrow Y$ un morphisme \deleted{loc. de \unsure{t}} propre de préschémas \inserted{loc} \unsure{noeths|noeth}, avec $Y$ intègre, $y_0\in Y$,\tikz[remember picture,overlay]{\coordinate (lgtb); \draw ([xshift=-5pt,yshift=-5pt]lgta) -- ([xshift=-5pt]lgta |- lgtb);} $y_1$ \deleted{le pt gén. de $Y$. \inserted{\unsure{une générisation}}} \inserted{le pt gén de $Y$}. On suppose \deleted{que l'on est dans le \unsure{l'un des cas suivants|cas des suivants}} \deleted{\unsure{intég.}} \deleted{\unsure{$y_1$ est un pt gén}}%
% vertical rule beside the first three lines of the hypotheses

\tikz[remember picture,overlay]{\coordinate (lgtc);}\margin{?}\deleted{(\unsure{a})} (a) $Y$ est [géom.\textsuperscript{t}] unibranche en $y_0$, \deleted{et} tte composante irréductible de $X$ domine $Y$.\tikz[remember picture,overlay]{\coordinate (lgtd); \draw ([xshift=-5pt]lgtc) -- ([xshift=-5pt]lgtc |- lgtd);}

\deleted{(b) $f$ est universellt ouvert} \inserted{(b) les composantes irréductibles de} \deleted{\unsure{ses en}} $X_1$ sont de dim $\geq k+1$, et $X_1$ est connexe en dim $\geq k$.

Alors (les composantes irréductibles de $X_0$ sont de dim $\geq k+1$) et $X_0$ est \uline{connexe} en dim $\geq k$.

L'assertion entre parenthèses est bien connue par le théorème des 2 dimensions, et n'utilise pas l'hypothèse unibranche. Celle-ci sert uniquemt à appliquer le \uline{th. de connexion}, qui assure que $X_0$ est connexe.

% archive sheet number: 22
\scanpage{24}
\mspage{- 2 -}
\noindent\hfill n\textsuperscript{o}355

d'o\handwritten{\`}{u} $\operatorname{Ass}_B(F\handwritten{/}\underline p F)\handwritten{\subset}\operatorname{Ass}_B(E\handwritten{\otimes}_A F)$, ce qui démontre l'inclusion $(\ast)$.

On va maintenant supposer que $A$ est \uline{noethérien}, et démontrer l'inclusion opposée. On procèdera par étapes :

(i) $E$ \uline{est un module de type fini et} $\operatorname{Ass}_A(E)$ \uline{est réduit à un élément} $\underline p$.

D'après le th. 1 du \handwritten{\S}{} 1, il existe une suite de composition $(E_i)_{0\handwritten{\leq} i\handwritten{\leq} n}$ de $E$ telle que $E_i\handwritten{/}E_{i+1}$ soit isomorphe à $A\handwritten{/}\underline p_i$, o\handwritten{\`}{u} $\underline p_i$ est un idéal premier de $A$. D'après le th.2 du \handwritten{\S}{} 1, tous les $\underline p_i$ contiennent $\underline p$. Du fait que $F$ est $A$-plat, les $E_i\handwritten{\otimes}_A F$ forment une suite de composition de $E\handwritten{\otimes}_A F$, et $(E_i\handwritten{\otimes}_A F)\handwritten{/}(E_{i+1}\handwritten{\otimes}_A F)$ s'identifie à $A\handwritten{/}\underline p_i\handwritten{\otimes}_A F=F\handwritten{/}\underline p_i F$. En appliquant la prop.3, on en déduit :

\[
\operatorname{Ass}_B(E\handwritten{\otimes}_A F)\handwritten{\subset}\handwritten{\bigcup}\operatorname{Ass}_B(F\handwritten{/}\underline p_i F)\,.
\]

Soit donc $\underline P\handwritten{\in}\operatorname{Ass}_B(E\handwritten{\otimes}_A F)\handwritten{\cap}\operatorname{Ass}_B(F\handwritten{/}\underline p_i F)$. Si $\underline p_i\handwritten{\ne}\underline p$, il existe $a\handwritten{\in}\underline p_i$, $a\handwritten{\notin}\underline p$. D'après le cor.2 à la prop.2 du \handwritten{\S}{} 1, l'homothétie de rapport $a$ est injective dans $E$, donc celle de rapport $\handwritten{\pi}(a)$ est injective dans $E\handwritten{\otimes}_A F$ ; mais d'autre part, $\handwritten{\pi}(a)$ annule $F\handwritten{/}\underline p_i F$, d'o\handwritten{\`}{u} $\handwritten{\pi}(a)\handwritten{\in}\underline P$, contradiction. Ainsi $\underline p_i$ est nécessairement égal à $\underline p$, ce qui démontre la proposition dans le cas considéré.

(ii) $E$ \uline{est un module de type fini}.

Soient $\underline p_i$ les éléments de $\operatorname{Ass}_A(E)$, et soit $0=\handwritten{\bigcap} Q_i$ une décomposition primaire réduite correspondante (cf. \handwritten{\S}{} 2, n\textsuperscript{o}2). Le module $E$ se plonge dans la somme directe des $E_i=E\handwritten{/}Q_i$, et $E\handwritten{\otimes}_A F$ se plonge dans la somme directe des $E_i\handwritten{\otimes}_A F$. On en déduit :

\[
\operatorname{Ass}_B(E\handwritten{\otimes}_A F)\handwritten{\subset}\handwritten{\bigcup}\operatorname{Ass}_B(E_i\handwritten{\otimes}_A F)\,.
\]

D'après (i), on a $\operatorname{Ass}_B(E_i\handwritten{\otimes}_A F)=\operatorname{Ass}_B(F\handwritten{/}\underline p_i F)$, d'o\handwritten{\`}{u} la proposition dans ce cas.

(iii) \uline{Cas général}.

% archive sheet number: 23
\scanpage{25}

Le lemme 4 nous ramène alors à vérifier le

\uline{Lemme 5} (\unsure{à énoncer|évidemment} avant le th).

\tikz[remember picture,overlay]{\coordinate (lgfa);}Sous les conditions du th, mais sans \deleted{\unsure{savoir ni a priori que|savoir a priori que}} \inserted{supposer $\deleted{\unsure{x}}y_0$ \deleted{géom.} \uline{unibranche}} en $y_0$, \inserted{ni $f$ propre,} pour tt pt fermé $x$ de $X_0$, \deleted{l'anneau local} $\operatorname{Spec}(\underline{\mathcal O}_{X_0,x})$ est connexe en dim $\geq k$\tikz[remember picture,overlay]{\coordinate (lgfb); \draw ([xshift=-5pt,yshift=4pt]lgfa) -- ([xshift=-5pt]lgfa |- lgfb);}
% slanted
\margin{dire plutôt : tout $\bar x$ tel que $\dim\bar x\leq k$, $\operatorname{Spec}\underline{\mathcal O}_{X_0,x})'$ est connexe}

En effet, le lemme 1 \inserted{+ \unsure{théor.|théorie} de dim} nous ramène facilement, utilisant l'hypothèse sur $X_1$, au cas où $X_1$ est \uline{irréductible} \deleted{de dim $\unsure{>|\geq} k+1$}, donc $X$ \uline{irréductible} dominant $Y$, \inserted{donc $=$ intègre comme $Y$.} On peut supposer par une réduction facile (?) au cas où ``\deleted{la formule des \inserted{\deleted{\unsure{en dépendant de type fini|essentielt de type fini}}} dimensions est vraie pour $\underline{\mathcal O}_{Y,y_0}$}\inserted{$\underline{\mathcal O}_{Y,y_0}$ est formellt équidimensionnel}''.
\margin{expliciter}

Donc

\[
\dim\underline{\mathcal O}_{X,x}=\underbrace{\dim\underline{\mathcal O}_{Y,y_0}}_{n}+\underbrace{\dim X_1}_{\geq k+1}
\]

% archive sheet number: 24
\scanpage{26}
\noindent\hfill ALGEBRE\quad COMMUTATIVE\hfill Rédaction n\textsuperscript{o}355

CHAPITRE IV - DECOMPOSITION PRIMAIRE

\uline{Décomposition primaire et extension des scalaires} (Etat 2)

\begin{center}-:-:-\end{center}

L'état 1 porte le n\textsuperscript{o}352. Après l'avoir lu (et même avant), Grothendieck a signalé au rédacteur qu'on pouvait généraliser le résultat principal, en remplaçant l'un des anneaux par un module. La démonstration s'étend sans changement, et c'est bien plus joli. Le rédacteur en a profité pour enlever quelques hypothèses noethériennes inutiles.

\begin{center}-:-:-\end{center}

Dans tout ce qui suit, on désigne par $A$ et $B$ deux anneaux commutatifs, et l'on considère un homomorphisme $\handwritten{\pi}:A\to B$. Cet homomorphisme fait de $B$ une $A$-algèbre. On se donne en outre un $B$-module $F$ que l'on suppose \uline{plat sur} $A$.

PROPOSITION 1.- \uline{Soit} $\underline p$ \uline{un idéal premier de} $A$. \uline{Pour tout} $\underline P\handwritten{\in}\operatorname{Ass}_B(F\handwritten{/}\underline p F)$ \uline{on a} $\handwritten{\pi}^{-1}(\underline P)=\underline p$.

Puisque $\handwritten{\pi}(\underline p)$ annule $F\handwritten{/}\underline p F$, on a $\handwritten{\pi}(\underline p)\handwritten{\subset}\underline P$, d'o\handwritten{\`}{u} $\underline p\handwritten{\subset}\handwritten{\pi}^{-1}(\underline P)$. D'autre part, si $a\handwritten{\in} A-\underline p$, l'homothétie de rapport $a$ est injective dans $A\handwritten{/}\underline p$, donc aussi (puisque $F$ est $A$-plat) dans $F\handwritten{/}\underline p F$, ce qui montre que $\handwritten{\pi}(a)\handwritten{\notin}\underline P$, CQFD.

PROPOSITION 2.- \uline{Soit} $E$ \uline{un} $A$-\uline{module}. \uline{On a} :

\[
\operatorname{Ass}_B(E\handwritten{\otimes}_A F)\handwritten{\supset}\handwritten{\bigcup}_{\underline p\handwritten{\in}\operatorname{Ass}_A(E)}\operatorname{Ass}_B(F\handwritten{/}\underline p F)
\tag*{\makebox[0pt][r]{\makebox[\linewidth][l]{$(\ast)$}}}
\]

\uline{Lorsque} $A$ \uline{est noethérien}, \uline{l'inclusion} $(\ast)$ \uline{est une égalité}.

Soit $\underline p\handwritten{\in}\operatorname{Ass}_A(E)$. Par définition de $\operatorname{Ass}_A(E)$, il existe une suite exacte :

\[
0\to A\handwritten{/}\underline p\to E\,.
\]

Puisque $F$ est plat sur $A$, on en déduit une suite exacte :

\[
0\to F\handwritten{/}\underline p F\to E\handwritten{\otimes}_A F\,,
\]

% archive sheet number: 25
\scanpage{27}

Donc divisant $\underline{\mathcal O}_{X,x}$ par un idéal défini par $\deleted{\unsure{k}}n$ éléments, on \deleted{\unsure{obtient|de ce fait}} trouve par le th général \deleted{\unsure{un}} (\unsure{appliquer|appliqué} : $\underline{\mathcal O}_{\unsure{X|x},x}$ intègre formellt équidim. de dim $\geq K+1=(n+k)+1$] un anneau qui est connexe en dim $\geq\deleted{(n+k)}\inserted{K-n}=k$, cqfd.

\uline{Cor. 1} au th. 2 Variante avec \deleted{composantes} ``\unsure{géométriquemt|géométriqt} \deleted{\unsure{int.|iréd.}} unibranche'', \unsure{*|``} \deleted{\unsure{et}} fibres géométriques.

\uline{Cor. 2} Variante avec $f$ propre univ. ouvert.

\uline{Cor. 3} $X\subset P^r_k$, \deleted{$X$ \unsure{de di}} \inserted{les composantes} de $\unsure{\alpha|X}$ de dim $\geq k+1$, et $X$ connexe en dim $\geq k$. Alors pour tt système de $\ell$ hyperplans $H_1\text{--}H_\ell$, \inserted{\unsure{le schéma|la section}} \deleted{\unsure{$Z$}} $X\cap H_1\cap-\cap H_\ell=Y$ \unsure{satisfait} les conditions analogues avec \unsure{au|$\ell$ au} lieu de $k$.
% slanted
\margin{Généraliser au cas d'un morphisme $X\longrightarrow P_k^r$.}
% slanted
\margin{Peut-on généraliser de façon analogue le th. local ? Facile, en utilisant le thé. \unsure{tel} \unsure{en raisonnant|correspondant} \unsure{relatif}.}

% archive sheet number: 26
\scanpage{28}

\deleted{\unsure{Théorie|théorie}.} \inserted{\uline{\unsure{Couples|Couple} de}} \uline{Lefschetz-Grauert}

\[
(\unsure{x|X},Y)
\]

% archive sheet number: 27
\scanpage{30}

\uline{\uline{III} \deleted{\unsure{1}}0}

\uline{Introduction}

Dans le par. 4, nous avons vu \deleted{que} lorsque $X$ est propre sur $Y$ \deleted{\unsure{un anneau local complet}} \inserted{\deleted{$I$-adique $A$}} \inserted{loc. noeth,} \deleted{et si on \unsure{connaît|considère}} et si $Y'$ est une partie fermée de $Y$, $X'=f^{-1}(Y')$, \deleted{la} \inserted{une propriété de} commutation des foncteurs \deleted{$R^i f$ et $R^i\hat f$ avec} ``images directes supérieures'' aux foncteurs complétion \unsure{Au long|(le long|le long} de $Y'$, resp $X'$]. Lorsque $Y$ était le spectre d'un anneau \deleted{local} noeth $I$-adique complet, et $Y'=V(A/I)$, ces résultats, complétés par un th d'existence (par 5), s'expriment \deleted{en disant} \inserted{en disant} \deleted{par} \inserted{que $F'\mathbin{-}F'_{/Y'}$ est une} équivalence de la catégorie des \deleted{\unsure{faisceaux}} \inserted{Modules} coh. sur $X'$ et de la catégorie des pro-Modules coh. sur $X'_{/Y'}$, \uline{compatible} avec la formation des modules de cohomologie.

Dans ce par, nous \deleted{généralisons} \inserted{montrons} que ces résultats peuvent se généraliser à certaines situations où $f:X'\to X$ n'est plus supposé propre [p. ex. $f$ une imm. ouverte \unsure{---|…}] \deleted{\unsure{ici}} la seule différence essentielle est que \deleted{\unsure{pr}} ici on \unsure{va trouver}
% slanted
\inserted{\unsure{dans l'énoncé des|des raisons ou} conditions \unsure{qui|que} pour des raisons \unsure{à voir|sont}, \unsure{par|pour} \unsure{nécessité|v. commentaire}.}
\deleted{qui s'expriment \unsure{d'une manière à préciser. Ceci|De manière à préciser ceci} \unsure{la vérité|le résultat} \unsure{essentiel est encore fourni} \unsure{un Th.}}
La méthode sera \inserted{\unsure{d'ailleurs toute parallèle à}} celle des par \deleted{\unsure{5|4}}4 et 5, \inserted{(cf N\textsuperscript{o}2)} en utilisant \deleted{des} \inserted{un} \uline{th. de finitude} convenable \inserted{(N\textsuperscript{o}1)} \deleted{\unsure{Ces de Ces derniers font}} \deleted{\unsure{faire intervenir}} faisant intervenir des propriétés \deleted{s'exprime :} \deleted{\unsure{intervenir reposent sur|intervenir reposant sur}} locales des faisceaux envisagés, qui s'expriment :

% archive sheet number: 29
\scanpage{31}

l'aide de la notion de \uline{profondeur}.

Les résultats \inserted{ainsi} obtenus sont l'outil technique principal pour un \deleted{\unsure{étude|dévelop}} \inserted{développement} systématique de \deleted{propriétés} \inserted{certains théorèmes} du type de ``Lefschetz'', comparant les propriétés géométriques d'un schéma $X$ à celles de certains \inserted{sous-schémas} diviseurs $Y$ \deleted{(sections hyperplanes)} \deleted{\unsure{en}} \inserted{sur $X$}, \uline{via} les propriétés du complété formel $X_{/Y}$. \deleted{\unsure{du}} \deleted{Ils \unsure{précisent|prennent} \unsure{une|en} \unsure{remarque|nouvelle} l'observation \unsure{faite|fait} \unsure{par} \unsure{application|appliqué}} Ils sont applicables par exemple \inserted{(\unsure{après des procédés connus|après du procédé général} de passage aux cônes \unsure{projetants|projetés})} au cas d'une variété projective $X$ et une section hyperplane convenable, moyennant des hypothèses \deleted{de reg\unsure{u}} convenables sur la profondeur de $X$. \deleted{\unsure{Ce qui précède} \unsure{Ceci est}}

Dans ce cas \deleted{\unsure{assez voisin du} à préciser} \deleted{la théorie} \deleted{\unsure{le reste}} \deleted{a été} esquissé par H. Grauert qui avait d'ailleurs, \deleted{\unsure{une}} (\unsure{observé|associé|Bourbaki} \deleted{\unsure{à}} \inserted{dès} 1958, (utilisant les th. de \deleted{\unsure{co}} nullité de Serre $[\,]$)$]$ \deleted{de Grauert que dans la situation précédente \unsure{et}} \deleted{\unsure{moyennant des hyp|la situation au-dessus}} \deleted{\unsure{supplémentaires}} \deleted{\unsure{de}}

(i) \inserted{Pour tt $F$ loc. libre,} $H^i(X,F)\to H^i(\hat X,\hat F)$ est bijectif pour des \deleted{petits} $i$ $\leq n$ $[n$ \deleted{dépendant de} \inserted{$\leq\operatorname{prof}F-2$}$]$

(ii) \unsure{d'une|d'un} part \unsure{si} un faisceau formel loc libre $\mathcal F$ sur $X_{/Y}$, \deleted{$\hat{\mathcal F}$ \unsure{est} \unsure{isom}} \inserted{et si $\operatorname{prof}(\mathcal O_{\hat X})\geq 2$,} alors pour \unsure{$n$|un} \deleted{coupe de sections ``} \unsure{grand|formel}, $\hat{\mathcal F}\unsure{(n)}$ est \unsure{engendré|engendre} par ses sections [$\unsure{\text{sauf}|\exists n}$] sur un ens fini -- $[$et \deleted{\unsure{à preuve le verrez|ce, sous les autres}} \inserted{en fait} \unsure{pratiquement|pendant} \deleted{\unsure{je}} comme \unsure{nous montrerons|une immersion}$]$, \deleted{Alors} \deleted{\unsure{des conditions|des Grothendieck}} \deleted{\unsure{C'est cette une des}} (Cela qui impliquait \inserted{p.ex.} déjà le th de Lefschetz pour le groupe fondamental énoncé en $[\,]$.)\margin{La théorie développée ici a pour but de préciser et de \deleted{\unsure{f}} généraliser les résultats de Grauert.}
% margin note written sideways

% archive sheet number: 30
\scanpage{32}

\begin{samepage}
\tikz[remember picture,baseline=(lgstart.base)]{\node[inner sep=0pt](lgstart){1.};} Généralisation des th. de finitude formelle, et des th. de comparaison. Application : un th. d'existence

2. Le Théorème de finitude locale et ses applications.

\deleted{3. Fermés de Lefschetz-Grauert\unsure{:|.|,} Remarques\unsure{.|,}}

\deleted{4. \unsure{Premières|Propriétés} propriétés \unsure{des} \uline{Exemples} \unsure{fondamentaux|fondament}} \inserted{1}

3. Cas particulier 1 : \deleted{\unsure{Etude|Étude}} \inserted{théorie} \unsure{des|d'en} \deleted{\unsure{spa}} schémas locaux.

4. Cas particulier 2 : \deleted{\unsure{Etude|Étude}} \inserted{théorie} des schémas \inserted{algébriques} projectifs.

5. Fermés de Lefschetz-Grauert : Remarques, exemples.

% slanted
\margin{\tikz[remember picture,baseline=(lgnote.base)]{\node[inner sep=1pt,text width=1.9cm,draw,align=left](lgnote){utilise IV, \S\ sur les morphismes étales....};}}
6. Propriétés liées aux $\pi_0$, \tikz[remember picture,baseline=(lgpi.base)]{\node[draw,ellipse,inner sep=1pt](lgpi){$\pi_1$};}, Pic.

[7. Propriétés de permanence]

\tikz[remember picture,baseline=(lgeight.base)]{\node[draw,ellipse,inner sep=1pt](lgeight){8.};} Théorèmes de pureté, et théorèmes de Lefschetz pour le $\pi_1$.

\deleted{\textcircled{\unsure{9|1}}} \tikz[remember picture,baseline=(lgnine.base)]{\node[inner sep=0pt](lgnine){9};} \uline{Anneaux} \deleted{\unsure{locaux}} \uline{quasi-factoriels} et \uline{th. de Lefschetz pour Pic}.

\tikz[remember picture,overlay]{\draw[->] (lgnote.east) to[bend left=12] (lgpi.north west); \draw[->] (lgnote.south east) to[bend right=12] (lgeight.west); \draw ([xshift=-4pt]lgstart.north west) -- ([xshift=-12pt]lgstart.north west) -- ([xshift=-12pt]lgnine.south west) -- ([xshift=-4pt]lgnine.south west); \draw ([xshift=-12pt,yshift=7pt]lgnine.north west) -- ([xshift=-5pt,yshift=7pt]lgnine.north west);}
\end{samepage}

% archive sheet number: 31
\scanpage{33}

\uline{Fermés de Lefschetz-Grauert}

1) Définitions, remarques

\tikz[remember picture,overlay]{\coordinate (lgda);}\uline{Déf} Soit $X$ un préschéma loc noeth, $Y$ une partie fermée, $n$ un entier. On dit que $Y$ a la propriété $(LG)_n$ dans $X$ \deleted{\unsure{jusqu'en|presque} \unsure{dim|si et|dans} \unsure{$n$|un}}, si pour tt $F$ loc. libre sur $X$,

\[
H^i(X,F)\mathbin{\unsure{-|\to}}H^i(X_{/Y},F_{/Y})\deleted{\unsure{\mathrm{est}|\mathrm{isom}}}
\tag*{\makebox[0pt][r]{\makebox[\linewidth][l]{$(\ast)$}}}
\]

est bijectif pour $i\leq n$, \inserted{et injectif pour $i=n+1$.}\margin{inutile}

On dit que $Y$ a la propriété $LG$ \inserted{faible dans $X$} si \deleted{elle a la propriété $(LG)_0$,} \inserted{pour} tout $F$ loc. libre sur $X$,

\[
H^0(X,F)\overset{\sim}{\longrightarrow}H^0(X_{/Y},F_{/Y})
\]

\tikz[remember picture,overlay]{\coordinate (lgdb); \draw ([xshift=-6pt,yshift=-3pt]lgda) -- ([xshift=-6pt]lgda |- lgdb); }[$i.e.$ $F\leadsto F_{/Y}$ est ple. fid. des Mod. loc. libres sur $X$ en ceux sur $X_{/Y}$]

On dit que $Y$ a la propriété $LG$ \deleted{\unsure{stricte|de stricte}} \inserted{dans $X$} si tt \unsure{pour|$V$} voisinage ouvert $U$ de $Y$, $Y$ \deleted{ait} \inserted{a la} propriété $LG$ \inserted{faible dans $U$}. \uline{\unsure{Enfin}} $Y$ a la propriété $LG$, \deleted{\unsure{stricte}} \inserted{\unsure{définitive|effective}} si elle a la propriété $LG$, \inserted{et si de plus} \deleted{pour}
\margin{$\Rightarrow$ Si $z\in X$ est tel \deleted{\unsure{fermé et ne ren-|qu'il se ren-}} \deleted{\unsure{contre pas $Y$|contre pas}} que $\bar z\cap Y=\varnothing$, alors $z\notin\operatorname{Ass}\underline{\mathcal O}_X$.}
tt $\mathcal F$ loc libre sur $X_{/Y}$ \deleted{est} est isom : un faisceau de la forme $F_{/Y}$, où $F$ loc libre sur un voisinage ouvert $U$ de $Y$.
% slanted
\margin{Question : $(\unsure{\mathrm{Lef}|\mathrm{LGf}|\mathrm{LG}})$ \deleted{\unsure{\#}} implique-t-elle \deleted{\unsure{formellement}} $\forall x\in X$ tq $\bar x\cap Y=\varnothing$, on a $\operatorname{prof}\mathcal O_{X,x}\geq 2$ ???}

Donc la propriété \textbf{LG} stricte signifie que les foncteurs $\mathcal L(U)\mathbin{-}\mathcal L(\unsure{W|V})$ $[\deleted{\unsure{\forall}}U\supset V\supset Y]$ sont pleinement fidèles, et

\[
\underset{U\supset Y}{\underset{\longrightarrow}{\tikz[remember picture,baseline=(lglim.base)]{\node[inner sep=0pt](lglim){$\operatorname{Lim}$};}}}\;
\underbrace{\mathcal L(U)}_{\substack{\text{Catégorie des}\\\text{Modules loc}\\\text{libres sur }U}}
\longrightarrow\mathcal L(X_{/Y})
\]
\tikz[remember picture,baseline=(lgcaption.base)]{\node[align=center,inner sep=1pt](lgcaption){``Pseudolim'' de\\catégories};}
\tikz[remember picture,overlay]{\draw[<->] (lgcaption.north east) -- (lglim.south west);}

est une équivalence de catégories.
% slanted
\margin{\uline{Exemples tout de suite}}

% archive sheet number: 32
\scanpage{34}

\uline{Remarque 1} La \uline{propriété $LG$ faible} \uline{équivaut} à $(LG)_0$.

En effet, il suffit de prouver qu'elle implique que $H^{\unsure{1}}(X,F)\to H^1(X_{/Y},F_{/Y})$ est injectif, i.e. que tt ext. de $\underline{\mathcal O}_X$ par $F$ qui splitte formellement, splitte régulièrement, ce qui est immédiat. [Question : Si $\unsure{\alpha|\ast}$ est un isom pour $i\leq n$, tt $F$ loc libre, est-il vrai que $\unsure{\alpha|\ast}$ est un mono pour \deleted{\unsure{tout}} $i=n+1$, tt $F$ loc libre ?.] Oui, prendre \inserted{une \deleted{\unsure{résolution}} \inserted{\unsure{résol. univers.|résolution}} \unsure{singulière|injective|régulière|injection}, $0\to F\to I\to G\to0$, et appliquer le foncteur $_{/Y}$, et le lemme des cinq]}

\uline{Remarque 2} Soit $(L_\alpha)$ une famille de Modules loc. libres sur $X$, telle que tt module loc. libre sur $X$ soit \inserted{$\simeq$} quotient d'une somme directe de Modules $\check L_\alpha$, i.e. $\simeq$ au \unsure{ann. dual|ann. d'un|ss. mod.|sous-module} \inserted{loc. facteur direct} d'une somme directe de Modules $L_\alpha$ $[$Ex : $1^\circ$ $X$ quasi \deleted{projectif} \inserted{affine}, $(L_\alpha)$ réduit : $\underline{\mathcal O}_X$ $2^\circ$ $X$ quasi-projectif, $(L_\alpha)$ formé des $\underline{\mathcal O}(m)$, $m\geq m_0$ $]$.

Alors $(LG)_n$ équivaut à la condition

\[
H^i(X,L_\alpha)\overset{\sim}{\longrightarrow}H^i(X_{/Y},L_{\alpha/Y})
\quad\left|\begin{aligned}&\text{isom pour }i\leq n\\[-2pt]&\text{inj pour }i=n+1\end{aligned}\right.
\]
% slanted
\margin{N.B. Le même argument marche si on n'exige pas la condition $H^{n+1}(\unsure{F|X})\to H^{n+1}(\hat{\unsure{F|X}})$ \uline{injectif}.}

\deleted{\unsure{Cet énoncé|Aut. rédac}} \deleted{[N.B. \unsure{Conclusion :|Contrairement :} \unsure{pour la|pourquoi} preuve \unsure{on mettra} une condition sur $i=n+1$ !]}. Dém par récurrence sur $n$, trivial si $n<-1$, ds $n+1<0$, supposons \deleted{\unsure{un}} que c'est prouvé

% archive sheet number: 33
\scanpage{35}

pour les $n'<n$. \deleted{Prouvons que} Nous savons donc déjà que \deleted{\unsure{$\ast$}}$(\ast^{\unsure{i}})$ est bijectif pour $i\leq n-1$, injectif pour $i=n$, on veut gagner un cran.

On a une suite exacte

\[
0\mathbin{\unsure{-|\to}}F\mathbin{-}L\mathbin{-}G\to0
\]

où \inserted{$G$ loc libre,} $L$ est une somme directe de modules loc. libres \deleted{\unsure{sommes|sur un}} $\simeq L_\alpha$, donc les assertions : pour $n$ sont vraies pour $L$.

Considérons l'hom de suites exactes

\diagram{\begin{tikzcd}[ampersand replacement=\&,column sep=small]
H^i(L) \arrow[r,dash] \arrow[d,"\wr\;\alpha"] \& H^i(G) \arrow[r,dash] \arrow[d,"\beta"] \& H^{i+1}(F) \arrow[r,dash] \arrow[d,"\varphi"] \& H^{i+1}(L) \arrow[r,dash] \arrow[d,"\wr\;\alpha'"] \& H^{i+1}(G) \arrow[d,"\beta'"] \\
H^i(\hat L) \arrow[r,dash] \& H^i(\hat G) \arrow[r,dash] \& H^{i+1}(\hat F) \arrow[r,dash] \& H^{i+1}(\hat L) \arrow[r,dash] \& H^{i+1}(\hat{\unsure{C|G}})
\end{tikzcd}}

pour $i=n-1,n$ donc $i+1=n,n+1$. Dans le premier cas, on veut prouver que la flèche médiane est surjective, dans le deuxième qu'elle est injective.

a) $i=n-1$, on sait par hyp. de réc. que $\beta$ est surjective, \deleted{(\unsure{est|ou|tout} bijective)} \deleted{il en est de même de $\alpha$, \unsure{enfin|alors} par l'hyp. de réc.} \inserted{de plus $\alpha'$ est surj par hyp. et $\beta'$ est injective.} Par lemme des 5, $\varphi$ est \deleted{bijective} surjective \inserted{(\deleted{\unsure{même|donc}} bijective)}

b) $\unsure{v|i}=n$, On sait par \deleted{a) l'hyp} \deleted{a) que $\beta$ est} \inserted{a)} \unsure{précemment|ci-dessus} que $\alpha$ et $\beta$ sont bijectives, ou $\alpha$ est surj, $\beta$ inj, de plus $\alpha'$ est inj.

% archive sheet number: 34
